% latexkit gallery — every component with the code that makes it.
% Build: latexkit gallery   (or: latexmk gallery.tex, after `latexkit gallery` wrote catalog.tex)
\documentclass[10pt,a4paper]{article}
\usepackage[lang=vi, all, parskip, margin=1.8cm]{latexkit}
\graphicspath{{../templates/_common/plots/}}

% Code on the left/top, result below
\NewTCBListing{demo}{!O{}}{enhanced, breakable, colback=white, colframe=black!15, boxrule=0.5pt,
  arc=2pt, listing above text, listing engine=listings,
  listing options={language={[LaTeX]TeX}, basicstyle=\ttfamily\footnotesize, numbers=none,
    keywordstyle=\color{accent}, moretexcs={gradbar,gradrow,gradcolumn,gradspark,gradheat,gradstack,midarrow,lkpipeline,lktimeline,lkdiagram,lktikzfig,
    lkfig,best,second,cmark,xmark,pmark,lkhl,lkheat,lkbar,lkcodefile,lkgroup,lkbibliography}},
  colupper=black, segmentation style={black!15}, #1}

% One OpenTikZ asset in the catalog: {type}{id}{name}
\NewTColorBox{lkasset}{m m m}{enhanced, colback=white, colframe=black!12, boxrule=0.4pt,
  arc=2pt, halign=center, valign=center, left=3pt, right=3pt,
  title={\scriptsize\sffamily\bfseries #2 \hfill\mdseries #1},
  coltitle=black!75, colbacktitle=black!4, toptitle=0.5pt, bottomtitle=0.5pt,
  after upper={\par\vspace{2pt}{\scriptsize\sffamily\color{black!60} #3}}}

\title{\sffamily\bfseries latexkit — Gallery}
\author{\sffamily Thành phần dùng chung cho ghi chú, paper, báo cáo, sách và slide}
\date{\sffamily\today}

\begin{document}
\maketitle
\tableofcontents

\section{Cài đặt trong tài liệu}
\begin{lkterminal}
\usepackage[lang=vi, table, diagram, plot, code, box, math, algo, science]{latexkit}
\usepackage[lang=en, all]{latexkit}          % every module, English labels
\usepackage[style=plain, accent=8B0000]{latexkit}  % classic headings, red accent
\end{lkterminal}

% ===================================================================
\section{Bảng (\texttt{table})}

\begin{demo}
\begin{lktable}[H]{Kết quả chính}{tab:demo}{l ccc}
  Phương pháp & PSNR\up & LPIPS\down & Real time \\ \midrule
  NeRF        & 26.1          & 0.31          & \xmark \\
  3DGS        & \second{27.4} & \second{0.22} & \cmark \\
  \lkhl Ours  & \best{28.8}   & \best{0.18}   & \cmark \\
\end{lktable}
\end{demo}

\begin{demo}
\begin{lktablex}[H]{Bảng giãn hết chiều ngang, ô màu theo giá trị}{tab:heat}{l C C C}
  Tiêu chí & A & B & C \\ \midrule
  Chất lượng & \lkheat{0.6} & \lkheat{0.9} & \lkheat{0.4} \\
  Tốc độ     & \lkbar{0.8}  & \lkbar{0.5}  & \lkbar{0.9} \\
  LiDAR      & \cmark & \pmark & \xmark \\
\end{lktablex}
\end{demo}

\IfFileExists{siunitx.sty}{\begin{demo}
\begin{center}
\begin{lktblr}{colspec={l S[table-format=2.1] S[table-format=1.3]}}
  {Method} & {PSNR} & {SSIM} \\
  Baseline & 25.1 & 0.810 \\
  Ours     & 28.8 & 0.893 \\
\end{lktblr}
\end{center}
\end{demo}

}{\textit{(Cần siunitx: texlive-science.)}}

\subsection{Đồ họa dữ liệu trong bảng (gradbars)}
\begin{demo}
\begin{gradtable}[width=\linewidth, header align=c]
  \gradcolumn{method}{Phương pháp}
  \gradcolumn[type=bar, mark=both, options={max=100, precision=1}]{score}{Điểm}
  \gradcolumn[type=number, mark=best, options={better=lower, precision=1}]{time}{Thời gian / ms}
  \gradrow{NeRF, 72.4, 18.2}
  \gradrow{3DGS, 92.7, 11.5}
  \gradrow{Ours, 88.1, 9.6}
\end{gradtable}

Trong câu: \gradbar[max=100, width=25mm]{72.4} \quad \gradspark{3,5,2,8,6}
\quad \gradheat[min=0, max=1]{0.7} \quad \gradstack{30,50,20}
\end{demo}

% ===================================================================
\section{Sơ đồ (\texttt{diagram})}

\subsection{Pipeline và timeline nhanh}
\begin{demo}
\lkpipeline{Dữ liệu thật, Tái dựng, Mô phỏng, Đánh giá}

\lkpipeline*[node distance=5mm]{Bước 1, Bước 2}
\end{demo}

\begin{demo}
\lktimeline{2020/NeRF, 2022/Instant-NGP, 2023/3DGS, 2024/{Street Gaussians}}
\end{demo}

\subsection{Các style node và cạnh}
\begin{demo}
\begin{tikzpicture}[lk]
  \node[terminal]                 (s) {Bắt đầu};
  \node[io,        right=of s]    (i) {Đọc ảnh};
  \node[decision,  right=of i]    (d) {Hợp lệ?};
  \node[process,   right=of d]    (p) {Xử lý};
  \node[datastore, below=of p]    (db) {Lưu};
  \node[note,      below=of i]    (n) {ghi chú};
  \draw[flow] (s) -- (i);  \draw[flow] (i) -- (d);
  \draw[flow] (d) -- node[elabel, above] {Có} (p);
  \draw[flow] (p) -- (db);
  \draw[loop] (d.south) |- ++(0,-0.9) -| node[elabel, pos=0.25] {Không} (i.south);
  \lkgroup{(d)(p)}{Lõi}
\end{tikzpicture}
\end{demo}

\begin{demo}
\begin{tikzpicture}[lk]
  \node[layer=9cm]                                 (a) {Ứng dụng};
  \node[layer=9cm, below=3mm of a]                 (b) {Mô hình};
  \node[layer=9cm, below=3mm of b, lkcolor=lkteal] (c) {Dữ liệu};
  \node[tag, anchor=east] at (a.east) {API};
\end{tikzpicture}
\end{demo}

\subsection{Cây phân loại (forest)}
\begin{demo}
\begin{forest}
  for tree={draw=accent, fill=accent!8, rounded corners=2pt, font=\sffamily\small,
            edge={draw=black!55}, l sep=6mm, s sep=3mm}
  [Biểu diễn 3D [Tường minh [Mesh] [Gaussians]] [Ẩn [NeRF] [SDF]]]
\end{forest}
\end{demo}

\subsection{Sơ đồ tuần tự (pgf-umlsd)}
\begin{demo}
\begin{sequencediagram}
  \newthread{c}{Client}
  \newinst[2]{s}{Server}
  \begin{call}{c}{request()}{s}{response}\end{call}
\end{sequencediagram}
\end{demo}

\subsection{Sơ đồ từ file: Mermaid, D2, Graphviz, draw.io}
File nằm trong \lkfile{diagrams/}; latexmk tự chuyển sang PDF.
\begin{demo}
\lkdiagram[max height=7cm]{diagrams/pipeline}   % pipeline.mmd

\lkdiagram{diagrams/architecture}                % architecture.d2

\lkdiagram[max width=0.7\linewidth]{diagrams/graph}  % graph.dot

\lkdiagram[max width=0.6\linewidth]{diagrams/sketch} % sketch.drawio
\end{demo}

\subsection{Hình OpenTikZ}
\begin{demo}
\lktikzfig[max height=6cm]{../assets/vendor/opentikz/templates/encoder-decoder/template}
\end{demo}
Thêm vào project: \texttt{latexkit asset add encoder-decoder} rồi \verb|\lktikzfig{figures/encoder-decoder}|.
Toàn bộ thư viện ở \cref{sec:catalog}.

\subsection{Kiến trúc mạng nơ-ron (PlotNeuralNet)}
\begin{demo}
\begin{tikzpicture}
  \pic[shift={(0,0,0)}] at (0,0,0) {RightBandedBox={name=c1, caption=conv1,
       xlabel={{"64","64"}}, zlabel=224, fill=\ConvColor, bandfill=\ConvReluColor,
       height=36, width={2,2}, depth=36}};
  \pic[shift={(0,0,0)}] at (c1-east) {Box={name=p1, fill=\PoolColor, opacity=0.5,
       height=28, width=1, depth=28}};
  \pic[shift={(2,0,0)}] at (p1-east) {RightBandedBox={name=c2, caption=conv2,
       xlabel={{"128","128"}}, zlabel=112, fill=\ConvColor, bandfill=\ConvReluColor,
       height=28, width={3,3}, depth=28}};
  \pic[shift={(2,0,0)}] at (c2-east) {Box={name=fc, caption=fc, xlabel={{"10",}},
       fill=\SoftmaxColor, height=3, width=1.5, depth=20}};
  \draw[connection] (p1-east) -- node{\midarrow} (c2-west);
  \draw[connection] (c2-east) -- node{\midarrow} (fc-west);
\end{tikzpicture}
\end{demo}
Hình có sẵn: \texttt{latexkit asset list nn} (AlexNet, VGG16, U-Net, FCN, HED...).

% ===================================================================
\section{Đồ thị (\texttt{plot})}
\begin{demo}
\begin{tikzpicture}
  \begin{axis}[lkplot, width=0.75\linewidth, height=5cm, xlabel={Bước}, ylabel={PSNR (dB)},
               legend pos=south east]
    \addplot table[x=step, y=baseline, col sep=comma] {../templates/_common/plots/results.csv};
    \addplot table[x=step, y=ours, col sep=comma]     {../templates/_common/plots/results.csv};
    \legend{Baseline, Ours}
  \end{axis}
\end{tikzpicture}
\end{demo}

\begin{demo}
\begin{tikzpicture}
  \begin{axis}[lkbar, symbolic x coords={NeRF, 3DGS, Ours}, ylabel={FPS}, ymin=0,
               width=0.6\linewidth, height=4.5cm]
    \addplot coordinates {(NeRF, 0.1) (3DGS, 120) (Ours, 95)};
  \end{axis}
\end{tikzpicture}
\end{demo}

Đồ thị bằng Python dùng style SciencePlots có sẵn trong \lkfile{plots/styles}:
\begin{lkcode}[Python]{plots/make\_plots.py}
from latexkit_plot import use_style, save
use_style("science", "ieee")      # hoặc "nature", "grid", "no-latex", ...
fig, ax = plt.subplots()
ax.plot(step, psnr)
save(fig, "psnr")                 # -> figures/psnr.pdf, chèn bằng \lkfig{psnr}{...}{...}
\end{lkcode}

\subsection{Phong cách figures4papers}
Bảng màu và kiểu trình bày lấy từ \href{https://github.com/ChenLiu-1996/figures4papers}{figures4papers}
(Chen Liu, CC BY-NC 4.0): xanh đậm cho phương pháp đề xuất, xanh lá cho cải tiến, đỏ hồng cho
baseline, chỉ giữ trục trái và trục dưới. Dùng được cho cả pgfplots (\texttt{lkf4p},
\texttt{lkf4pbar}, màu \texttt{f4pblue}, \texttt{f4pgreenA..C}, \texttt{f4predA..B}, \texttt{f4pred},
\texttt{f4pgray}) và matplotlib (\lkfile{plots/f4p.py}).
\begin{demo}
\begin{tikzpicture}
  \begin{axis}[lkf4pbar, symbolic x coords={AUROC, AUPRC, F1}, ylabel={Điểm}, ymin=0.3, ymax=0.95,
               width=0.75\linewidth, height=5cm, bar width=8pt, legend columns=3,
               legend style={at={(0.5,1.03)}, anchor=south}]
    \addplot coordinates {(AUROC, 0.71) (AUPRC, 0.38) (F1, 0.33)};
    \addplot coordinates {(AUROC, 0.77) (AUPRC, 0.52) (F1, 0.40)};
    \addplot[fill=f4pblue, draw=black] coordinates {(AUROC, 0.84) (AUPRC, 0.63) (F1, 0.49)};
    \legend{Baseline A, Baseline B, Ours}
  \end{axis}
\end{tikzpicture}
\end{demo}

Bằng Python: \texttt{latexkit asset add f4p-bars-comparison} chép script mẫu vào \lkfile{plots/},
sửa số liệu ở đầu file rồi chạy \texttt{latexkit plot}. Có 12 mẫu (\texttt{latexkit asset list plot}),
xem ở \cref{sec:catalog}.
\begin{lkcode}[Python]{plots/ket\_qua.py}
from f4p import PALETTE, LARGE, apply_publication_style, make_grouped_bar, finalize_figure
apply_publication_style(LARGE)            # font 24, trục dày 3: hình so sánh lớn
fig, ax = plt.subplots(figsize=(12, 5))
make_grouped_bar(ax, ["A", "B"], [[0.78, 0.80], [0.92, 0.88]], ["Baseline", "Ours"],
                 colors=[PALETTE["red_1"], PALETTE["blue_main"]], hatches=["//", ""], annotate=True)
finalize_figure(fig, "ket_qua")           # -> figures/ket_qua.pdf
\end{lkcode}

% ===================================================================
\section{Code (\texttt{code})}
\begin{lkcode}[Python][label=lst:train]{Vòng lặp huấn luyện}
for epoch in range(num_epochs):     # chú thích tiếng Việt vẫn hiển thị đúng
    loss = criterion(model(x), y)
    loss.backward()
\end{lkcode}
Viết là \verb|\begin{lkcode}[Python][label=lst:train]{Tiêu đề}|; tham chiếu bằng \cref{lst:train}.
Lấy code từ file: \verb|\lkcodefile[Python]{code/train.py}|. Lệnh terminal:
\begin{lkterminal}
$ latexkit new report bao-cao -t "Báo cáo" -m all
\end{lkterminal}

% ===================================================================
\section{Hộp ghi chú (\texttt{box})}
\begin{demo}
\begin{keyidea} Ý chính của phần. \end{keyidea}
\begin{info} Thông tin thêm. \end{info}
\begin{tip} Mẹo. \end{tip}
\begin{warning} Lưu ý. \end{warning}
\begin{question}[Câu hỏi khi đọc] Paper giải quyết vấn đề gì? \end{question}
\begin{examplebox} Ví dụ. \end{examplebox}
\begin{summary} Tóm tắt. \end{summary}
\end{demo}

% ===================================================================
\section{Toán và thuật toán (\texttt{math}, \texttt{algo}, \texttt{science})}
\IfFileExists{siunitx.sty}{\begin{demo}
\begin{definition} Hàm $f:\R^n\to\R$ lồi nếu $f(\lambda x+(1-\lambda)y)\le\lambda f(x)+(1-\lambda)f(y)$. \end{definition}
$\hat\theta = \argmin_\theta \norm{\mat{A}\vect{x}-\vect{b}}^2$, \quad \SI{9.81}{\metre\per\second\squared}, \quad \num{12345.678}
\end{demo}
}{\textit{(Cần siunitx: texlive-science.)}}

\begin{demo}
\begin{algorithm}[H]
  \caption{Gradient descent}
  \begin{algorithmic}[1]
    \Require $\theta_0, \eta$
    \While{chưa hội tụ}
      \State $\theta \gets \theta - \eta\nabla\mathcal{L}(\theta)$ \Comment{cập nhật}
    \EndWhile
  \end{algorithmic}
\end{algorithm}
\end{demo}

% ===================================================================
\clearpage
\section{Thư viện hình: OpenTikZ, PlotNeuralNet và figures4papers}\label{sec:catalog}
Mỗi ô là một file \texttt{standalone} trong \lkfile{assets/vendor/opentikz} (hình theo CC0 1.0,
code theo MIT) hoặc \lkfile{assets/vendor/plotneuralnet} (MIT, loại \texttt{nn}), biên dịch cùng latexkit,
hoặc một script matplotlib trong \lkfile{assets/vendor/figures4papers} (CC BY-NC 4.0: ghi nguồn, không dùng
cho mục đích thương mại; loại \texttt{plot}, số liệu chỉ là ví dụ).
Thêm vào project bằng \texttt{latexkit asset add <id>}.

\begin{tcbraster}[raster columns=3, raster equal height=rows, raster column skip=6pt, raster row skip=6pt]
\InputIfFileExists{catalog.tex}{}{\textit{Chạy \texttt{latexkit gallery} để tạo danh mục.}}
\end{tcbraster}

\end{document}
